\section*{अध्याय \ref{ch:b1:diffeq} --- रैखिक अवकल समीकरण}

\begin{solution}{exo:b1:diffeq:1}
समांगी: $y_h = \lambda\,\eu^{-2x}$। विशिष्ट: $y_p = c\,\eu^{3x}$ आज़माइए
($3$ $r + 2$ का मूल नहीं है): $3c + 2c = 1$, $c = \frac15$। व्यापक हल: $y =
\frac{\eu^{3x}}{5} + \lambda\,\eu^{-2x}$। $y(0) = 1$ के साथ: $\frac15 +
\lambda = 1$, $\lambda = \frac45$, अतः $y = \frac{\eu^{3x} +
4\,\eu^{-2x}}{5}$।
\end{solution}

\begin{solution}{exo:b1:diffeq:2}
$\intoo{0}{+\infty}$ पर $x$ से भाग दीजिए: $y' - \frac{2}{x}\,y = x^2$। यहाँ
$A(x) = -2\ln x$, $\eu^{-A(x)} = x^2$: समांगी हल $\lambda x^2$। \omterm{thm:b1:diffeq:voc}{अचरों का
विचरण}: $\mu'(x) = x^2 \cdot x^{-2} = 1$, अतः $\mu = x + \lambda$ और
\[
y(x) = x^3 + \lambda x^2, \qquad \lambda \in \R .
\]
\emph{जाँच:} $x(3x^2 + 2\lambda x) - 2(x^3 + \lambda x^2) = x^3$।
\end{solution}

\begin{solution}{exo:b1:diffeq:3}
$y'' - 3y' + 2y = 0$: मूल $1$ और $2$; $\;y = \lambda\,\eu^{x} +
\mu\,\eu^{2x}$।

$y'' + 4y' + 4y = 0$: द्विक मूल $-2$; $\;y = (\lambda + \mu x)\,\eu^{-2x}$।

$y'' - 2y' + 5y = 0$: मूल $1 \pm 2\iu$; $\;y = \eu^{x}(\lambda\cos 2x +
\mu\sin 2x)$।
\end{solution}

\begin{solution}{exo:b1:diffeq:4}
समांगी: मूल $\pm 1$, $y_h = \lambda\,\eu^x + \mu\,\eu^{-x}$। बहुपदीय दाएँ
पक्ष के साथ विशिष्ट ($\gamma = 0$ मूल नहीं है): $y_p = ax^2 + bx + c$;
प्रतिस्थापित करने पर $2a - (ax^2 + bx + c) = x^2$ से $a = -1$, $b = 0$, $c =
2a = -2$ मिलते हैं: $y_p = -x^2 - 2$। व्यापक हल $y = -x^2 - 2 + \lambda\eu^x
+ \mu\eu^{-x}$।

कोशी: $y(0) = -2 + \lambda + \mu = 0$ और $y'(0) = \lambda - \mu = 1$:
$\lambda = \frac32$, $\mu = \frac12$। अतः $y = -x^2 - 2 + \frac{3\eu^x +
\eu^{-x}}{2}$।
\end{solution}

\begin{solution}{exo:b1:diffeq:5}
$a(x) = \tan x$, $A(x) = -\ln(\cos x)$ (वैध: अंतराल पर $\cos > 0$),
$\eu^{-A} = \cos x$: समांगी हल $\lambda\cos x$। \omterm{thm:b1:diffeq:voc}{अचरों का विचरण}: $\mu'(x) =
\sin 2x \cdot \frac{1}{\cos x} = 2\sin x$, अतः $\mu = -2\cos x + \lambda$ और
\[
y(x) = -2\cos^2 x + \lambda \cos x .
\]
\emph{जाँच:} $y' = 4\cos x \sin x - \lambda\sin x$ और $y\tan x = -2\cos
x\sin x + \lambda \sin x$; उनका योग $2\cos x\sin x = \sin 2x$ है, जैसा
अपेक्षित था।
\end{solution}

\begin{solution}{exo:b1:diffeq:6}
$\chi(r) = r^2 - 4r + 3 = (r-1)(r-3)$: $\gamma = 1$ सरल मूल है ($m = 1$)।
$y_p = x(ax + b)\,\eu^x$ आज़माइए। $u = ax^2 + bx$ के साथ,
\[
y_p'' - 4y_p' + 3y_p = \bigl(u'' + (2 - 4)u' + \chi(1) u\bigr)\eu^x
= \bigl(2a - 2(2ax + b)\bigr)\eu^x .
\]
$(2x + 1)\eu^x$ से पहचान कीजिए: $-4a = 2$ और $2a - 2b = 1$, अतः $a =
-\frac12$, $b = -1$। व्यापक हल:
\[
y = -\Bigl(\frac{x^2}{2} + x\Bigr)\eu^{x} + \lambda\,\eu^{x} +
\mu\,\eu^{3x}, \qquad (\lambda, \mu) \in \R^2 .
\]
\end{solution}

\begin{solution}{exo:b1:diffeq:7}
समांगी: $y_h = \lambda\cos 2x + \mu\sin 2x$।

दायाँ पक्ष $x$ ($\gamma = 0$ मूल नहीं है): $4(ax + b) = x$ के साथ $y_1 = ax
+ b$: $y_1 = \frac x4$।

दायाँ पक्ष $\sin 2x = \Im(\eu^{2\iu x})$, जहाँ $2\iu$ $r^2 + 4$ का सरल मूल
है: $z = c\,x\,\eu^{2\iu x}$ आज़माइए; तब $z'' + 4z = 4\iu c\,\eu^{2\iu x}$,
जो $c = \frac{1}{4\iu} = -\frac{\iu}{4}$ के लिए $\eu^{2\iu x}$ के बराबर है।
अतः $z = -\frac{\iu x}{4}(\cos 2x + \iu \sin 2x)$ और $y_2 = \Im(z) =
-\frac{x\cos 2x}{4}$।

अध्यारोपण से:
\[
y = \frac{x}{4} - \frac{x\cos 2x}{4} + \lambda\cos 2x + \mu\sin 2x .
\]
\end{solution}

\begin{solution}{exo:b1:diffeq:8}
समीकरण $T' + kT = 20k$ का अचर विशिष्ट हल $20$ है और समांगी हल
$\lambda\eu^{-kt}$ हैं: $T(t) = 20 + \lambda\,\eu^{-kt}$, और $T(0) = 80$ से
$\lambda = 60$ मिलता है:
\[
T(t) = 20 + 60\,\eu^{-kt} .
\]
$T(10) = 50$: $\eu^{-10k} = \frac12$, अतः $k = \frac{\ln 2}{10}$। फिर $T(t)
= 25$ के लिए $\eu^{-kt} = \frac{5}{60} = \frac{1}{12}$ आवश्यक है, अर्थात्
\[
t = \frac{\ln 12}{k} = 10\,\frac{\ln 12}{\ln 2} \approx 35.8
\text{ मिनट।}
\]
\end{solution}

\begin{solution}{exo:b1:diffeq:9}
\begin{enumerate}
  \item $\chi(r) = r^2 + 2\varepsilon r + 1$, $\Delta = 4(\varepsilon^2 -
  1)$। $\varepsilon \in \intco{0}{1}$ के लिए: मूल $-\varepsilon \pm
  \iu\sqrt{1 - \varepsilon^2}$, अतः $\omega = \sqrt{1 - \varepsilon^2}$ के
  साथ $y = \eu^{-\varepsilon t}\bigl(\lambda\cos\omega t + \mu\sin\omega
  t\bigr)$। $\varepsilon = 1$ के लिए: द्विक मूल $-1$, $y = (\lambda + \mu
  t)\,\eu^{-t}$। $\varepsilon > 1$ के लिए: वास्तविक मूल $r_\pm =
  -\varepsilon \pm \sqrt{\varepsilon^2 - 1}$, दोनों $< 0$, और $y =
  \lambda\eu^{r_+t} + \mu\eu^{r_-t}$।
  \item $\varepsilon \in \intoo{0}{1}$ के लिए: $\abs y \leq
  \eu^{-\varepsilon t}(\abs\lambda + \abs\mu) \to 0$। $\varepsilon = 1$ के
  लिए: $(\lambda + \mu t)\eu^{-t} \to 0$ (चरघातांकी बहुपद को हरा देता है,
  \cref{prop:b1:functions:powerrules})। $\varepsilon > 1$ के लिए: दोनों
  चरघातांकी क्षय होते हैं, क्योंकि $r_\pm < 0$ (वस्तुतः $\sqrt{\varepsilon^2
  - 1} < \varepsilon$)। $\varepsilon = 0$ के लिए: $y = \lambda\cos t +
  \mu\sin t$ का आयाम अचर $\sqrt{\lambda^2 + \mu^2} \neq 0$ रहता है, जब तक $y
  = 0$ न हो।
  \item $R = \sqrt{\lambda^2 + \mu^2} > 0$ के साथ $\lambda\cos\omega t +
  \mu\sin\omega t = R\cos(\omega t - \varphi)$ लिखिए। $y$ के शून्य वही हैं
  जो $\cos(\omega t - \varphi)$ के हैं (गुणनखंड $\eu^{-\varepsilon t}$ कभी
  लुप्त नहीं होता): $\omega t - \varphi \equiv \frac\pi2 \pmod \pi$, जो
  अंतराल $\frac{\pi}{\omega} = \frac{\pi}{\sqrt{1 - \varepsilon^2}}$ वाली
  समांतर श्रेणी है।
\end{enumerate}
\end{solution}

\begin{solution}{exo:b1:diffeq:10}
$x = y = 0$ रखिए: $2f(0) = 2f(0)^2$, और $f(0) \neq 0$ से $f(0) = 1$ मिलता
है। $x$ नियत कीजिए और समीकरण का $y$ के सापेक्ष दो बार अवकलन कीजिए:
\[
f''(x+y) + f''(x-y) = 2 f(x) f''(y) .
\]
$y = 0$ रखने पर: $\;2f''(x) = 2 f(x) f''(0)$, अर्थात्
\[
f''(x) = c\,f(x), \qquad c = f''(0).
\]
\emph{स्थिति $c = \omega^2 > 0$:} $f(x) = \lambda\cosh\omega x +
\mu\sinh\omega x$; $f(0) = 1$ से $\lambda = 1$ मिलता है। फलनीय समीकरण में
रखने पर और योग-सूत्रों (\cref{prop:b1:functions:hyprules}) का प्रयोग करने पर
समीकरण $\mu = 0$ पर बाध्य कर देता है ($\sinh\omega x \sinh\omega y$ के
गुणांकों की तुलना कीजिए या $x = y$ पर मान लीजिए): $f = \cosh\omega x$, जो
$\cosh(x+y) + \cosh(x-y) = 2\cosh x\cosh y$ को सचमुच संतुष्ट करता है।

\emph{स्थिति $c = -\omega^2 < 0$:} इसी प्रकार $f(x) = \cos\omega x$ ($\omega
\neq 0$), जो समीकरण को संतुष्ट करता है।

\emph{स्थिति $c = 0$:} $f(0) = 1$ के साथ $f$ ऐफ़ीन है: $f(x) = 1 + \mu x$;
समीकरण $\mu = 0$ पर बाध्य करता है, जो वर्जित है ($f$ अचर नहीं है)।

निष्कर्ष: हल $f(x) = \cos\omega x$ और $f(x) = \cosh\omega x$, $\omega > 0$
हैं।
\end{solution}

\begin{solution}{exo:b1:diffeq:11}
$z(t) = y(\eu^t)$ रखिए, जिससे $x > 0$ के लिए $y(x) = z(\ln x)$। तब
\[
y'(x) = \frac{z'(\ln x)}x,
\qquad
y''(x) = \frac{z''(\ln x) - z'(\ln x)}{x^2} ,
\]
और समीकरण में प्रतिस्थापित करने पर:
\[
x^2y'' - xy' + y = \bigl(z'' - z'\bigr) - z' + z
= z'' - 2z' + z = 0 .
\]
\omterm{def:b1:diffeq:linear2}{अभिलक्षणिक बहुपद} $(r - 1)^2$: द्विक मूल $1$, अतः $z(t) = (\lambda + \mu
t)\,\eu^t$, और चर $x = \eu^t$ में लौटकर:
\[
y(x) = (\lambda + \mu\ln x)\,x,
\qquad \lambda, \mu \in \R .
\]
\end{solution}

\begin{solution}{exo:b1:diffeq:12}
\begin{enumerate}
  \item प्रत्येक अंतराल पर मानकीकृत रूप $y' - \frac2x\,y = 0$ में: $A(x) =
  -2\ln\abs x$, अतः हल $\intoo0{+\infty}$ पर $y = a x^2$ और
  $\intoo{-\infty}0$ पर $y = b x^2$ हैं, जिनके अचर स्वतंत्र हैं
  (\cref{thm:b1:diffeq:homogeneous1})।
  \item मान लीजिए $x \geq 0$ के लिए $y = ax^2$ और $x < 0$ के लिए $bx^2$।
  प्रत्येक खुली अर्धरेखा पर $y$ अवकलनीय है और $xy' = 2y$। $0$ पर: अंतर-भागफल
  $\frac{y(h) - y(0)}h = ah$ या $bh$ $0$ की ओर जाते हैं, अतः $y'(0) = 0$
  विद्यमान है, और $x = 0$ पर समीकरण $0 \cdot y'(0) = 2y(0) = 0$ कहता है: जो
  संतुष्ट है। अतः $y$ पूरे $\R$ पर समीकरण को हल करता है।
  \item $\R$ पर हल ठीक यही जोड़े हुए फलन हैं: प्रथम कोटि के समीकरण के लिए
  \emph{दो}-प्राचल वाला कुल। \cref{thm:b1:diffeq:voc} से कोई विरोध नहीं है,
  क्योंकि यहाँ उसकी परिकल्पनाएँ विफल हैं: $y' + a(x)y = 0$ के रूप में लिखने
  पर गुणांक $a(x) = -\frac2x$ $0$ पर संतत नहीं है --- वस्तुतः परिभाषित ही
  नहीं है --- अतः $\R$ ऐसा अंतराल नहीं है जिस पर वह प्रमेय लागू हो। $0$ की
  विचित्रता दोनों अर्धरेखाओं को अलग कर देती है, और मान $y(0) = 0$ अनिवार्य
  हो जाता है, जो एक ओर से दूसरी ओर कोई सूचना नहीं ले जाता। $x_0 \neq 0$ पर
  हर कोशी आँकड़ा हल को केवल उसी अर्धरेखा पर निर्धारित करता है जिसमें $x_0$
  है।
\end{enumerate}
\end{solution}

\begin{solution}{pb:b1:diffeq:1}
\textbf{1.} $0 < \lambda < \omega_0$ के लिए $\chi(r) = r^2 + 2\lambda r +
\omega_0^2$, $\Delta = 4(\lambda^2 - \omega_0^2) < 0$: मूल $-\lambda \pm
\iu\omega_d$, अतः \cref{thm:b1:diffeq:homogeneous2} से
\[
x(t) = \eu^{-\lambda t}\bigl(\lambda_1\cos\omega_d t +
\mu_1\sin\omega_d t\bigr), \qquad (\lambda_1, \mu_1) \in \R^2 .
\]
$\lambda = 0$ के लिए: $x = \lambda_1\cos\omega_0 t + \mu_1\sin\omega_0 t$।
क्रांतिक ($\lambda = \omega_0$) और अति-अवमंदित ($\lambda > \omega_0$)
अवस्थाएँ वही हैं जो \cref{exo:b1:diffeq:9} की हैं (समय के पुनर्मापन के बाद):
$(\lambda_1 + \mu_1 t)\eu^{-\lambda t}$, अथवा $r_\pm = -\lambda \pm
\sqrt{\lambda^2 - \omega_0^2}$ के साथ $\eu^{r_\pm t}$ के संयोजन।

\textbf{2.} अल्प-अवमंदित: $\abs x \leq \eu^{-\lambda t}(\abs{\lambda_1} +
\abs{\mu_1}) \to 0$। क्रांतिक: $(\lambda_1 + \mu_1 t)\eu^{-\lambda t} \to
0$, क्योंकि चरघातांकी बहुपदों को हरा देते हैं
(\cref{prop:b1:functions:powerrules})। अति-अवमंदित: $r_- < r_+ = -\lambda +
\sqrt{\lambda^2 - \omega_0^2} < 0$, क्योंकि $\sqrt{\lambda^2 - \omega_0^2} <
\lambda$; दोनों चरघातांकी क्षय होते हैं।

\textbf{3.} $x'' = -2\lambda x' - \omega_0^2 x$ का प्रयोग करते हुए $(H)$ के
किसी हल पर:
\[
\mathcal E'(t) = x'x'' + \omega_0^2 x x'
= x'\bigl(-2\lambda x' - \omega_0^2 x\bigr) + \omega_0^2 xx'
= -2\lambda\,x'^2 \leq 0 .
\]
यदि $x(t_0) = x'(t_0) = 0$, तो $\mathcal E(t_0) = 0$; $\mathcal E$ ऋणेतर और
अवर्धमान है, अतः $\intco{t_0}{+\infty}$ पर $\mathcal E \equiv 0$, जिससे वहाँ
$x \equiv 0$ अनिवार्य हो जाता है; $t \leq t_0$ के लिए वही तर्क $\tilde x(t)
= x(2t_0 - t)$ पर चलाइए, जो अवमंदन $-\lambda$ वाले समीकरण को हल करता है पर
उसमें भी $\tilde{\mathcal E}(t_0) = 0$ और $\tilde{\mathcal E} \geq 0$ के साथ
$\tilde{\mathcal E}' = +2\lambda\tilde x'^2 \geq 0$ है; ऋणेतर, अह्रासमान और
$\intoc{-\infty}{t_0}$ के दाएँ सिरे पर शून्य होने का अर्थ है सर्वत्र शून्य।
अतः $\R$ पर $x \equiv 0$ --- अद्वितीयता की एक ऊर्जा-उपपत्ति, जो सभी $\lambda
\geq 0$ के लिए वैध है।

\textbf{4.} $x(t + T_d) = R\,\eu^{-\lambda t}\eu^{-\lambda T_d}\cos(\omega_d
t + 2\pi - \varphi) = \eu^{-\lambda T_d}x(t)$। प्रति छद्म-आवर्तकाल
संकुचन-गुणक $\delta = \lambda T_d = \frac{2\pi\lambda}{\omega_d}$ के साथ
$\eu^{-\delta}$ है। $\omega_0 = 1$, $\lambda = 0.1$ के लिए: $\omega_d =
\sqrt{0.99} = 0.99499$, अतः $\delta = \frac{0.62832}{0.99499} = 0.6315$: हर
झूला अपने आयाम का $\eu^{-0.63} \approx 53\%$ बनाए रखता है।

\textbf{5.} आयाम-गुणक $\eu^{-\lambda t}$ है, जो $t = \frac1\lambda$ में
$\eu$ से घटता है। उस अंतराल में $\frac{1/\lambda}{T_d} =
\frac{\omega_d}{2\pi\lambda}$ छद्म-आवर्तकाल आते हैं। $\lambda \ll \omega_0$
के लिए $\omega_d \approx \omega_0$, और यह $\approx
\frac{\omega_0}{2\pi\lambda} = \frac Q\pi$ है। $Q = 300$ वाली गिटार की तार
लगभग सौ आवर्तकाल तक गूँजती है; $Q = 1$ वाला दरवाज़े का अवमंदक एक भी पूरा
नहीं करता।

\textbf{6.} बाएँ पक्ष में $z\,\eu^{\iu\Omega t}$ प्रतिस्थापित करने पर
$z\,(-\Omega^2 + 2\iu\lambda\Omega + \omega_0^2)\,\eu^{\iu\Omega t}$ मिलता
है, जो ठीक $z = \frac{A}{\omega_0^2 - \Omega^2 + 2\iu\lambda\Omega}$ के लिए
$A\,\eu^{\iu\Omega t}$ के बराबर है (हर अशून्य है: उसका काल्पनिक भाग
$2\lambda\Omega > 0$ है)। चूँकि गुणांक वास्तविक हैं, वास्तविक भाग $x_p =
\Re\bigl(z\eu^{\iu\Omega t}\bigr)$ दाएँ पक्ष $\Re\bigl(A\eu^{\iu\Omega
t}\bigr) = A\cos\Omega t$ वाले समीकरण को हल करता है।

\textbf{7.} $D = (\omega_0^2 - \Omega^2)^2 + 4\lambda^2\Omega^2$ और $\varphi
\in \intoo0\pi$ के साथ $\omega_0^2 - \Omega^2 + 2\iu\lambda\Omega =
\sqrt{D}\,\eu^{\iu\varphi}$ लिखिए (काल्पनिक भाग $2\lambda\Omega$ धनात्मक
है), अतः $\tan\varphi = \frac{2\lambda\Omega}{\omega_0^2 - \Omega^2}$। तब $z
= \frac{A}{\sqrt D}\eu^{-\iu\varphi}$ और
\[
x_p(t) = \Re\Bigl(\frac A{\sqrt D}\,\eu^{\iu(\Omega t -
\varphi)}\Bigr) = R\cos(\Omega t - \varphi),
\qquad R = \frac A{\sqrt D} .
\]

\textbf{8.} जब $\Omega \to 0^+$: $D \to \omega_0^4$, अतः $R \to
A/\omega_0^2$ और $\varphi \in \intoo0{\frac\pi2}$ के साथ $\tan\varphi \to
0^+$: $\varphi \to 0$। द्रव्यमान बल का अनुसरण अर्ध-स्थैतिक रूप से करता है,
और उसका विस्थापन बल/दृढ़ता होता है। जब $\Omega \to +\infty$: $D \sim
\Omega^4$, अतः $R \sim A/\Omega^2 \to 0$, और $\varphi \to \pi$ (सम्मिश्र
संख्या $\omega_0^2 - \Omega^2 + 2\iu\lambda\Omega$ कोणांक $\to \pi$ के साथ
द्वितीय चतुर्थांश में चली जाती है): द्रव्यमान मुश्किल से हिलता है, और
कला-विरोध में --- जड़त्व प्रभावी है।

\textbf{9.} \cref{thm:b1:diffeq:structure2}~(1) से प्रत्येक हल $x = x_p +
x_h$ है, जहाँ $x_h$ $(H)$ को हल करता है; प्रश्न 2 से $x_h(t) \to 0$, अतः
$x(t) - x_p(t) \to 0$: सभी हल उसी स्थायी अवस्था की ओर अभिसरित होते हैं।
आरंभिक शर्तें केवल क्षणिक भाग को आकार देती हैं।

\textbf{10.} यदि $x$ आवर्ती हल है, तो $x - x_p = x_h$ $(H)$ का ऐसा आवर्ती हल
है जो $+\infty$ पर $0$ की ओर जाता है; सीमा $0$ वाला आवर्ती फलन तत्समक रूप से
$0$ होता है (एक आवर्तकाल के उसके मान सदा दोहराते हैं, अतः हर मान $0$ की ओर
जाते किसी अनुक्रम की सीमा है)। अतः $x = x_p$।

\textbf{11.} यहाँ $\lambda = 1$, $\omega_0^2 = 2$, $\Omega = 1$, $A = 1$: $z
= \frac1{2 - 1 + 2\iu} = \frac{1 - 2\iu}5$, अतः
\[
x_p = \Re\Bigl(\frac{(1 - 2\iu)(\cos t + \iu\sin t)}5\Bigr)
= \frac{\cos t + 2\sin t}5 .
\]
समांगी: $r^2 + 2r + 2$ के मूल $-1 \pm \iu$ हैं: $x_h = \eu^{-t}(C\cos t +
S\sin t)$। शर्तें: $x(0) = \frac15 + C = 0$ से $C = -\frac15$ मिलता है;
अवकलन करने पर $x'(0) = \frac25 - C + S = 0$ से $S = C - \frac25 = -\frac35$
मिलता है। अतः
\[
x(t) = \underbrace{\frac{\cos t + 2\sin t}5}_{\text{स्थायी}}
- \underbrace{\eu^{-t}\,\frac{\cos t + 3\sin
t}5}_{\text{क्षणिक}} ,
\]
जहाँ क्षणिक भाग $\eu^{-t}$ की तरह मरता है।

\textbf{12.} प्रसार करने पर $u_r = \omega_0^2 - 2\lambda^2$ के साथ $g(u) =
u^2 - 2(\omega_0^2 - 2\lambda^2)u + \omega_0^4 = (u - u_r)^2 + g(u_r)$, और
\[
g(u_r) = \omega_0^4 - u_r^2 = (\omega_0^2 - u_r)(\omega_0^2 +
u_r) = 2\lambda^2\bigl(2\omega_0^2 - 2\lambda^2\bigr) =
4\lambda^2(\omega_0^2 - \lambda^2) .
\]
यदि $2\lambda^2 < \omega_0^2$, तो $u_r > 0$ स्वीकार्य वर्ग-आवृत्ति है: $g$
का वहाँ कठोर निम्नतम है, अतः $R = A/\sqrt g$ का $\Omega_r = \sqrt{u_r} =
\sqrt{\omega_0^2 - 2\lambda^2}$ पर कठोर अधिकतम है, जहाँ $R_{\max} =
A/\sqrt{g(u_r)} = \frac{A}{2\lambda\sqrt{\omega_0^2 - \lambda^2}}$।

\textbf{13.} $R(0) = A/\omega_0^2$, अतः
\[
\frac{R_{\max}}{R(0)}
= \frac{\omega_0^2}{2\lambda\sqrt{\omega_0^2 - \lambda^2}}
= \frac{\omega_0}{2\lambda}\cdot
\frac{\omega_0}{\sqrt{\omega_0^2 - \lambda^2}}
= Q\,\Bigl(1 - \frac{\lambda^2}{\omega_0^2}\Bigr)^{-1/2}
\geq Q .
\]
दुर्बल अवमंदन के लिए संशोधन-गुणक $1$ के निकट है: \omterm{ex:b1:diffeq:oscillation}{अनुनाद} स्थैतिक विस्थापन को
मूलतः $Q$ से गुणा कर देता है।

\textbf{14.} $u = \Omega^2$ के साथ $V(\Omega)^2 = \frac{A^2 u}{g(u)}$। उसके
अवकलज का चिह्न $g(u) - u\,g'(u) = (\omega_0^2 - u)^2 + 4\lambda^2 u -
u\bigl(2(u - \omega_0^2) + 4\lambda^2\bigr) = (\omega_0^2 - u)^2 +
2u(\omega_0^2 - u) = (\omega_0^2 - u)(\omega_0^2 + u)$ का चिह्न है, जो $u <
\omega_0^2$ के लिए धनात्मक और उससे आगे ऋणात्मक है: अतः प्रत्येक अवमंदन के
लिए कठोर अधिकतम ठीक $\Omega = \omega_0$ पर है। वहाँ $\varphi \in \intoo0\pi$
के साथ $\tan\varphi$ फट पड़ता है: $\varphi = \frac\pi2$, और $x_p'(t) =
-R\Omega\sin(\Omega t - \frac\pi2) = R\Omega\cos(\Omega t)$ बल के ठीक समान
कला में है: इष्टतम शक्ति-अंतरण।

\textbf{15.} $R = R_{\max}/\sqrt2 \iff g(u) = 2g(u_r) \iff (u - u_r)^2 =
g(u_r) = 4\lambda^2(\omega_0^2 - \lambda^2)$, जिससे
\[
u_\pm = u_r \pm 2\lambda\sqrt{\omega_0^2 - \lambda^2},
\qquad
\Omega_+^2 - \Omega_-^2 = 4\lambda\sqrt{\omega_0^2 - \lambda^2} .
\]
फिर $\Omega_+ - \Omega_- = \frac{\Omega_+^2 - \Omega_-^2}{\Omega_+ +
\Omega_-}$, और $\lambda \ll \omega_0$ के लिए दोनों $\Omega_\pm \approx
\omega_0$: $\Omega_+ - \Omega_- \approx \frac{4\lambda\omega_0}{2\omega_0} =
2\lambda$, अतः $\frac{\omega_0}{\Omega_+ - \Omega_-} \approx
\frac{\omega_0}{2\lambda} = Q$। अनुनाद-शिखर की चौड़ाई नापना उसका \omterm{pb:b1:diffeq:1}{गुणता कारक}
नापना है।

\textbf{16.} $Q = 10$; $\Omega_r = \sqrt{1 - 2(0.05)^2} = \sqrt{0.995} =
0.9975$; $R_{\max} = \frac1{2 \times 0.05 \sqrt{1 - 0.0025}} = \frac1{0.1
\times 0.99875} = 10.01$; स्थैतिक अनुक्रिया $R(0) = 1$; पट्ट-चौड़ाई $\approx
2\lambda = 0.1$। $1$ ऊँचाई के पठार पर $\approx Q$ ऊँचाई की एक ऊँची पतली नोक।

\textbf{17.} यदि $2\lambda^2 \geq \omega_0^2$, तो सभी $u > 0$ के लिए $u_r
\leq 0$ और $g'(u) = 2(u - u_r) > 0$: $\intoo0{+\infty}$ पर $g$ कठोरतः बढ़ता
है, अतः $R = A/\sqrt g$ $R(0) = A/\omega_0^2$ से कठोरतः घटता है: अनुक्रिया
शून्य आवृत्ति पर सबसे बड़ी है और कोई शिखर नहीं है।

\textbf{18.} $\gamma = \iu\Omega$ $r^2 + \omega_0^2$ का मूल नहीं है (क्योंकि
$\Omega \neq \omega_0$), अतः $m = 0$ के साथ \cref{met:b1:diffeq:particular}
$x_p = \frac{A\cos\Omega t}{\omega_0^2 - \Omega^2}$ देता है (प्रतिस्थापित
करके जाँचिए: कोज्या का $-\Omega^2 + \omega_0^2$ गुना)। व्यापक हल:
\[
x(t) = \frac{A\cos\Omega t}{\omega_0^2 - \Omega^2}
+ \lambda_1\cos\omega_0 t + \mu_1\sin\omega_0 t .
\]

\textbf{19.} $x(0) = 0$ $\lambda_1 = -\frac A{\omega_0^2 - \Omega^2}$ पर
बाध्य करता है, और $x'(0) = 0$ $\mu_1 = 0$ पर:
\[
x(t) = \frac{A}{\omega_0^2 - \Omega^2}\,
\bigl(\cos\Omega t - \cos\omega_0 t\bigr)
= \frac{2A}{\omega_0^2 - \Omega^2}\,
\sin\Bigl(\frac{(\omega_0 - \Omega)t}2\Bigr)
\sin\Bigl(\frac{(\omega_0 + \Omega)t}2\Bigr) ,
\]
यह गुणनफल-सूत्र $\cos a - \cos b = 2\sin\frac{b - a}2\sin\frac{b + a}2$ को
$a = \Omega t$, $b = \omega_0 t$ के साथ लगाने से मिलता है।

\textbf{20.} $\omega_0$ के निकट $\Omega$ के लिए दूसरी ज्या तीव्र आवृत्ति
$\frac{\omega_0 + \Omega}2 \approx \omega_0$ पर दोलन करती है, जबकि पहली
आवृत्ति $\frac{\abs{\omega_0 - \Omega}}2$ का मंद अन्वालोप है: तीव्र दोलन का
आयाम अन्वालोप-आवर्तकाल $\frac{2\pi}{\abs{\omega_0 - \Omega}}$ के साथ
बढ़ता-घटता है (प्रति अन्वालोप-आवर्तकाल दो \omterm{pb:b1:diffeq:1}{विस्पंद}) और अधिकतम
$\frac{2A}{\abs{\omega_0^2 - \Omega^2}}$ तक पहुँचता है। जब $\Omega \to
\omega_0$, तब \omterm{pb:b1:diffeq:1}{विस्पंद} और धीमे (आवर्तकाल $\to \infty$) तथा और ऊँचे (आयाम $\to
\infty$) दोनों हो जाते हैं।

\textbf{21.} $t$ नियत कीजिए। जब $\Omega \to \omega_0$:
\[
\frac{2A}{\omega_0^2 - \Omega^2}\sin\Bigl(\frac{(\omega_0 -
\Omega)t}2\Bigr)
= \frac{2A}{\omega_0 + \Omega}\cdot
\frac{\sin\bigl(\frac{(\omega_0 - \Omega)t}2\bigr)}
{\omega_0 - \Omega}
\longrightarrow \frac{2A}{2\omega_0}\cdot\frac t2
= \frac{At}{2\omega_0} ,
\]
जबकि $\sin\bigl(\frac{(\omega_0 + \Omega)t}2\bigr) \to \sin(\omega_0 t)$:
सीमा $x_\infty(t) = \frac{At\sin\omega_0 t}{2\omega_0}$ है। सीधी जाँच: $C =
\frac A{2\omega_0}$ के साथ $x_\infty = Ct\sin\omega_0 t$ में $x_\infty'' =
2C\omega_0\cos\omega_0 t - C\omega_0^2 t \sin\omega_0 t$ है, अतः $x_\infty''
+ \omega_0^2 x_\infty = 2C\omega_0\cos\omega_0 t = A\cos\omega_0 t$, जहाँ
$x_\infty(0) = 0$ और $x_\infty'(0) = C\sin 0 + C\omega_0 \cdot 0 \cdot \cos
0 = 0$ --- और $\omega_0 = A = 1$ के लिए यह ठीक
\cref{ex:b1:diffeq:oscillation} है। \omterm{ex:b1:diffeq:oscillation}{अनुनाद} विस्पंदों का अपकर्ष है: अन्वालोप
की पहली उभार, अनंत लंबाई तक खिंची हुई।

\textbf{22.} अवमंदन के बिना अनुनादी आयाम रैखिक रूप से और बिना परिबंध बढ़ता
है; अवमंदन $\lambda > 0$ के साथ यह वृद्धि $R_{\max} \approx Q\,\frac
A{\omega_0^2}$ पर संतृप्त हो जाती है। क्रियाविधि: क्षय ऊर्जा को ऐसी दर से
हटाता है जो आयाम के साथ बढ़ती है (प्रश्न 23), अतः संचय ठीक तब रुक जाता है जब
अवमंदन उतनी ही तेज़ी से ऊर्जा जलाने लगता है जितनी तेज़ी से प्रणोदन उसे देता
है।

\textbf{23.} $x_p' = -R\Omega\sin(\Omega t - \varphi)$ के साथ, एक आवर्तकाल
पर औसत $\langle\cos^2\rangle = \langle\sin^2\rangle = \frac12$ और
$\langle\sin\cos\rangle = 0$ देते हैं:
\[
\langle P_{\mathrm{diss}}\rangle
= 2\lambda\,R^2\Omega^2\,\langle\sin^2(\Omega t -
\varphi)\rangle = \lambda R^2\Omega^2 ;
\]
और $\sin(\Omega t - \varphi) = \sin\Omega t\cos\varphi - \cos\Omega
t\sin\varphi$ का प्रसार करने पर:
\[
\langle P_{\mathrm{in}}\rangle
= -AR\Omega\,\bigl\langle\cos\Omega t\,\sin(\Omega t -
\varphi)\bigr\rangle
= AR\Omega\,\frac{\sin\varphi}2 .
\]
चूँकि $\sin\varphi = \frac{2\lambda\Omega}{\sqrt D} = \frac{2\lambda\Omega
R}A$, यह $\frac{AR\Omega}2 \cdot \frac{2\lambda\Omega R}A = \lambda
R^2\Omega^2$ है: प्रविष्ट और क्षयित शक्तियाँ ठीक-ठीक संतुलित हैं --- यही
स्थायी अवस्था का परिभाषक गुणधर्म है।

\textbf{24.} (क) संरचना प्रमेय ने प्रत्येक हल को स्थायी अवस्था और क्षणिक भाग
में बाँट दिया (प्रश्न 9--11) और अद्वितीयता को समांगी समस्या तक सीमित कर
दिया। (ख) सम्मिश्र विधि ने विशिष्ट हल की खोज को सम्मिश्र संख्याओं के एक ही
भाग में बदल दिया (प्रश्न 6), जहाँ आयाम और कला एक \omterm{def:b1:complex:field}{मापांक} तथा एक कोणांक से पढ़
लिए जाते हैं। (ग) \omterm{ex:b1:diffeq:oscillation}{अनुनाद} वक्र शुद्ध फलन-अध्ययन है --- $u = \Omega^2$ में एक
द्विघात, उसका निम्नतम, उसके स्तर-समुच्चय --- \cref{ch:b1:functions} की शैली
में (प्रश्न 12--17)।

\textbf{25.} मुक्त और अवमंदित: क्षयशील छद्म-दोलन, जीवनकाल $\frac1\lambda$,
लगभग $\frac Q\pi$ झूले। प्रणोदित और अवमंदित: क्षणिक भाग मर जाते हैं और
प्रणोदन-आवृत्ति पर एक अद्वितीय ज्यावक्रीय स्थायी अवस्था बची रहती है, जिसका
आयाम $\omega_0$ के पास शिखर पर होता है (ऊँचाई $\approx Q \times$ स्थैतिक,
चौड़ाई $\approx \frac{\omega_0}Q$) और कला $0$ से $\omega_0$ पर $\frac\pi2$
से होती हुई $\pi$ तक बहती है। मुक्त और अनवमंदित: शाश्वत दोलन। प्रणोदित और
अनवमंदित: \omterm{pb:b1:diffeq:1}{विस्पंद}, जो यथार्थ अनुरणन पर रैखिक रूप से बढ़ते \omterm{ex:b1:diffeq:oscillation}{अनुनाद} में
अपकर्षित हो जाते हैं। एक ही विमाहीन संख्या $Q = \frac{\omega_0}{2\lambda}$
सब कुछ साध देती है --- शिखर की ऊँचाई, पट्ट-चौड़ाई और क्षणिक जीवनकाल उसी
घुंडी के तीन पाठ हैं। आगे की कड़ी: $x'' + 2\lambda x' + \omega_0^2x$ को
प्रथम कोटि के निकाय के रूप में फिर से लिखना \cref{ch:b1:matrices} तथा स्नातक
वर्ष 2 के खंड की आव्यूह विधियाँ खोल देता है, और किसी भी आवर्ती प्रणोदन का
ज्यावक्रीय फलनों में वियोजन (फूरिये श्रेणी, स्नातक वर्ष 3 का खंड) इस समस्या
के एक-आवृत्ति विश्लेषण को सार्वभौमिक निर्माण-खंड बना देता है: प्रत्येक
आवृत्ति के लिए हल कीजिए और अध्यारोपित कर दीजिए।
\end{solution}
